\documentclass[11pt]{article}
\usepackage[margin=1in]{geometry}
\usepackage{parskip}
\usepackage{xcolor}
\usepackage{titlesec}
\titleformat{\section}[block]{\large\bfseries\color{blue!50!black}}{}{0pt}{}
\titleformat{\subsection}[block]{\bfseries\color{blue!50!black}}{}{0pt}{}
\setlength{\parskip}{0.5em}

\title{Speaker script for \emph{Deep Learning for Mean Field Games with Non-Separable Hamiltonians}}
\author{Companion to \texttt{presentation\_mfdgm.tex}\\
        Hikmatullo Ismatov}
\date{}

\begin{document}
\maketitle

\textit{Each section below corresponds to one slide.  Read each section as one
``unit of speech'' --- roughly 60--90 seconds per slide.  Total: about 40--45
minutes.  Trim freely; the most important slides to deliver carefully are
4--7 (foundations), 13--16 (GANs), 18--19 (Hopf), 26--28 (MFDGM), and 33--37
(our work).}

% ================================================================
\section*{Slide 1 --- Title page}

Hello everyone. Thank you for being here today.

My name is Hikmatullo Ismatov, I am a graduate student at KAUST, and this talk
is joint work with my supervisor Professor Diogo Gomes.

Today I will walk us through a paper that was published in 2023 by
Mouhcine Assouline and Badr Missaoui, called \emph{Deep Learning for Mean
Field Games with Non-Separable Hamiltonians}. It appeared in the journal
\emph{Chaos, Solitons and Fractals}.

I will explain the paper slowly, defining every term, so that even if you
have never heard of mean field games or GANs, you can follow.

At the end of the talk I will also briefly present a new method we are
developing in our group, which builds directly on this paper.

% ================================================================
\section*{Slide 2 --- What this talk covers}

The plan of the talk has seven parts.

We will start with foundations: what is a Mean Field Game, and why is it hard
to solve in high dimensions. Then we will look at the two big families of
deep-learning solvers for Mean Field Games: GAN-based methods and physics-
informed methods.

Before reading the paper, we need to understand three background tools:
what a GAN is, what the Hopf formula is, and what APAC-Net is. These three
sit behind everything in the paper.

After that, we will read the paper itself, both the algorithm and the
convergence theory. We will look at the experiments at the end.

And finally, I will tell you what our group is doing on top of this paper.

Important: please interrupt me at any time with questions. I will define
every term before I use it.

% ================================================================
\section*{Slide 3 --- Part 1 separator}

\textit{(Section page, automatic.)}  Take a breath, drink water.  Say:
``Part 1: foundations.''

% ================================================================
\section*{Slide 4 --- What is a Mean Field Game?}

Imagine the following picture. A thousand drivers on the same highway. Each
driver wants to reach their destination quickly. But the road has limited
capacity, so where it is crowded, everyone slows down, and where it is empty,
everyone speeds up.

This means each driver's best strategy depends on what everyone else does.
And everyone else's strategy depends on each driver. The whole situation is
circular.

A Mean Field Game is the mathematical theory that resolves this circularity.
It was developed independently by Lasry and Lions in 2007, and by Huang,
Caines, and Malham\'e in 2006. The idea is to find a Nash equilibrium: a
strategy from which no single agent benefits by changing their mind alone.

The word ``mean field'' comes from physics. It means: no individual agent is
big enough to matter. Only the average behaviour of the entire crowd matters.

% ================================================================
\section*{Slide 5 --- The two quantities we care about}

At Nash equilibrium, there are two objects we need to compute.

The first is the population density, denoted by $\rho(t, x)$. This tells you,
at time $t$, how many drivers are at position $x$.

The second is the value function, denoted by $\varphi(t, x)$. This is a more
abstract object. Think of $\varphi$ as a map of regret. It tells each agent:
``if you are at position $x$ at time $t$ and you play optimally from now
on, how much cost will you accumulate?''

So $\rho$ is observable in the world. You can stand on a bridge and count
how many cars pass.  But $\varphi$ lives inside each driver's head. It is the
agent's internal optimal-cost function.

% ================================================================
\section*{Slide 6 --- The MFG system}

These two objects satisfy two coupled equations.

The first is called the Hamilton-Jacobi-Bellman equation, or HJB for short.
It governs the value function $\varphi$, and it runs backward in time. You
start from a terminal condition at time $T$ and integrate backward to time
zero. This is because the value function depends on what happens in the
future.

The second is the Fokker-Planck equation, or FP. It governs the density $\rho$
and runs forward in time. You start from the initial density at time zero
and integrate forward.

The two equations are coupled. Look at the diagram on the right. The gradient
of $\varphi$ tells each agent where to go: this becomes the drift in the
Fokker-Planck equation. And the density $\rho$ enters the cost in the HJB
equation. So neither equation can be solved alone.

The number $\nu$ is the diffusion coefficient. When $\nu$ is zero, the
problem is deterministic; when $\nu$ is positive, there is random noise.

% ================================================================
\section*{Slide 7 --- The Hamiltonian H}

A Hamiltonian sounds intimidating, but it is just a fancy name for a cost
function.

Specifically, $H$ is defined as the Legendre transform of the running cost
$L$. The formula on the slide says: $H$ of $x$, $\rho$, $p$ equals the
supremum over velocities $v$ of $-v\cdot p$ minus $L$. The variable $p$
plays the role of momentum.

You do not need to remember the formula. The key thing to remember is: $H$
is given by the modeller. It is data of the problem. It packages all the
agent's costs and interactions into one function.

For example, if the running cost is one half $v$ squared minus $f_0$ of
$\rho$, the Hamiltonian works out to one half $p$ squared minus $f_0$ of
$\rho$. The momentum part and the density part stay separate.

In the MFG, the variable $p$ becomes the gradient of $\varphi$. So $H$
takes input $x$, $\rho$, and the gradient of $\varphi$.

% ================================================================
\section*{Slide 8 --- Separable vs non-separable Hamiltonians}

This is the most important distinction in the entire talk. Please pay close
attention.

A Hamiltonian is called separable if it splits as $H_0$ of $x$, $p$ minus
$f_0$ of $x$, $\rho$. The momentum part and the density part appear in
different terms and are never multiplied together. An example is one half
$p$ squared minus lambda times $\rho$. Notice $\rho$ and $p$ sit in
separate terms.

A Hamiltonian is non-separable if $\rho$ and $p$ appear multiplied together
somewhere. The canonical example is the traffic flow model of Huang and
collaborators from 2019: one half $p$ squared minus the term one minus
$\rho$, multiplied by $p$. The factor one minus $\rho$ multiplies $p$
directly. Density and momentum are entangled.

Why does this matter? Because some methods, specifically the GAN family,
only work when $H$ is separable. As soon as $\rho$ and $p$ are entangled,
those methods break. The traffic flow model is the canonical hard example
that motivates the entire paper.

% ================================================================
\section*{Slide 9 --- Part 2 separator}

\textit{(Section page.)}  Say: ``Part 2: why is this hard?''

% ================================================================
\section*{Slide 10 --- The curse of dimensionality}

Why can't we just solve the MFG with classical numerical methods?

Classical methods, like finite differences, put down a grid of points
covering the domain. If you have a one-dimensional problem with one hundred
grid points, you have a hundred unknowns. Easy.

Two dimensions with a hundred points per axis: ten thousand unknowns. Still
fine.

Three dimensions: one million unknowns. Doable.

Ten dimensions: ten to the twentieth power. That is already very hard.

A hundred dimensions: ten to the two hundredth power. This is more than the
number of atoms in the entire universe.

This is called the curse of dimensionality. Classical methods die in high
dimensions.

But many MFG applications, like multi-agent reinforcement learning, live in
hundreds of dimensions. We need a different approach.

% ================================================================
\section*{Slide 11 --- Neural networks to the rescue}

A neural network is just a parameterised function. The parameters are
called weights. We can tune the weights.

There is a famous theorem by Hornik from 1991, the Universal Approximation
Theorem. It says: a neural network with a single hidden layer and enough
neurons can approximate any continuous function to any desired accuracy. And
this works in any dimension, without a grid.

So the strategy is: replace the unknowns $\varphi$ and $\rho$ with neural
networks $\varphi_\omega$ and $\rho_\theta$, where omega and theta are the
weights. Then optimise the weights so the networks solve the PDE.

Crucially, the number of parameters of a neural network does not blow up
exponentially with dimension. A hundred-dimensional network with a hundred
neurons per layer has only about ten thousand parameters. Tractable.

This is why deep learning is the right tool for high-dimensional MFGs.

% ================================================================
\section*{Slide 12 --- Two ways to train the networks}

There are two big families of methods.

On the left, the physics-informed family. Sometimes called PINN, sometimes
DGM. The idea is: plug the network into the PDE and minimise the residual.
You sample random collocation points, evaluate the network and its
derivatives there, plug into the PDE, and minimise the squared error.
This works for any Hamiltonian. Today's paper sits here.

On the right, the variational or GAN-based family. The idea is: rewrite the
PDE as a saddle-point problem, a min-max, and let two networks play a game.
Examples are APAC-Net by Lin and collaborators in 2021, and MFGAN by
Ruthotto and collaborators in 2020. These methods only work when $H$ is
separable, because they need the Hopf formula to turn the PDE into a
saddle-point.

Our work, which I will present at the end, bridges these two columns.

% ================================================================
\section*{Slide 13 --- Part 3 separator}

\textit{(Section page.)}  Say: ``Part 3: what is a GAN?''

% ================================================================
\section*{Slide 14 --- Generative models: the goal}

To understand GANs, forget about MFGs for a moment.

The setup is this: I show you a million photographs of cats. I ask you to
write a computer program that, when I press a button, outputs a new cat
photo that I have never seen, but it has to look real.

The mathematical version: you are given samples from an unknown probability
distribution, called $P_\text{data}$. You want to build a program that
produces fresh samples from that distribution.

This is the generative modelling problem. Generative Adversarial Networks,
or GANs, are one solution to it.

% ================================================================
\section*{Slide 15 --- GAN: a two-player game}

A GAN, introduced by Goodfellow and collaborators in 2014, consists of two
neural networks that play a game.

The first network is called the generator, denoted $G_\theta$. It takes
random noise $z$ as input and produces a fake sample $x$. The goal of the
generator is to make the fakes look indistinguishable from real samples.

The second network is called the discriminator, denoted $D_\omega$. It takes
a sample, real or fake, and outputs the probability that the sample is real.
The goal of the discriminator is to correctly catch the fakes.

The two networks compete. The objective on the slide says: the discriminator
wants to maximise the probability of correct classification, on real and
fake samples. The generator can only influence the second term, and tries
to minimise it, that is, to fool the discriminator.

At equilibrium of this game, if both networks are powerful enough, the
generator produces samples exactly from the true data distribution.

Notice: this is itself a Nash equilibrium, of a two-player game. That is the
first hint that GANs and MFGs are linked.

% ================================================================
\section*{Slide 16 --- The Wasserstein GAN}

The original GAN has practical problems. It can be unstable to train. So a
better version was introduced by Arjovsky, Chintala, and Bottou in 2017,
called the Wasserstein GAN.

In the Wasserstein GAN, instead of asking the discriminator real or fake,
we ask it to assign a score to every sample. The discriminator is restricted
to be 1-Lipschitz, meaning it cannot change too fast. The generator's job
is to make its samples score high.

There is a beautiful theorem from optimal transport theory, called the
Kantorovich-Rubinstein duality. It says: the maximum value of this game is
equal to the Wasserstein-1 distance between the real distribution and the
generator's distribution.

This means: training a Wasserstein GAN is literally minimising the
Wasserstein distance. The training loss has a clean geometric meaning at
every step. That is why W-GANs are easier to train.

% ================================================================
\section*{Slide 17 --- Why are GANs and MFGs linked?}

Now look at the two formulas on the slide.

The first is the Wasserstein GAN, which is a min-max problem over two
networks.

The second is the saddle-point form of a separable MFG, derived by Cao,
Guo, and Lauri\`ere in 2021. It is also a min-max problem, with $\rho$
playing the role of the generator and $\varphi$ playing the role of the
discriminator.

The two have exactly the same structure. A distribution-producing object
plays against a function being optimised. This shared structure is the
reason GAN-based methods can solve separable MFGs.

But, very important: deriving the saddle-point form of an MFG requires
the Hopf formula. And the Hopf formula requires separability. We will see
why on the next two slides.

% ================================================================
\section*{Slide 18 --- Part 4 separator}

\textit{(Section page.)}  Say: ``Part 4: the Hopf formula.''

% ================================================================
\section*{Slide 19 --- The Legendre transform}

The Legendre transform is an old idea from physics, used to convert between
position and momentum.

For a function $L$ of velocity $v$, the Legendre transform $L^*$ is defined
on the slide. It takes the supremum over $v$ of the inner product $p \cdot
v$ minus $L$ of $v$.

Two concrete examples on the slide:

Example one. If $L$ of $v$ equals one half $v$ squared, the Legendre
transform is one half $p$ squared.

Example two, more interesting. If $L$ depends on velocity and density, $L$
of $v$, $\rho$ equals one half $v$ squared plus $f$ of $\rho$, then the
Legendre transform in $v$ only is one half $p$ squared plus the same $f$ of
$\rho$.

Notice the crucial point: the $f$ of $\rho$ piece was constant in $v$, so
it sat there during the transform. The Legendre transform preserved the
separation between $\rho$ and $p$.

This is exactly why separability matters: the Legendre transform plays
nicely with $\rho$ when $\rho$ does not interact with $v$.

% ================================================================
\section*{Slide 20 --- The Hopf formula}

The Hopf formula is a way to write the solution of a Hamilton-Jacobi
equation as an optimisation problem.

For the simplest HJ equation, $\partial_t \varphi$ plus $H$ of gradient $\varphi$
equals zero, the Hopf formula says: $\varphi$ of $t$, $x$ equals the sup
over $y$ of the inf over $z$ of an expression involving the initial data,
the time, and the Hamiltonian. You can see it on the slide.

The right-hand side is a min-max problem.

This is why the Hopf formula is so important for our story. A min-max problem
is exactly what GAN-style training can attack. The Hopf formula is the
bridge that turns an MFG into a GAN.

% ================================================================
\section*{Slide 21 --- Why the Hopf formula needs separability}

For an MFG, the Hamilton-Jacobi equation has $H$ of $x$, $\rho$, gradient
$\varphi$. If $H$ is separable, then $f_0$ of $x$, $\rho$ is constant in
$p$. You can pull it out, treat it as a known source term for $\varphi$, and
apply the Hopf formula to the remaining $H_0$ of $x$, gradient $\varphi$.
This gives a clean min-max problem, exactly what APAC-Net uses.

If $H$ is non-separable, the density and momentum are entangled inside $H$.
The Legendre transform variable cannot factor cleanly. The Hopf formula does
not give a min-max reformulation any more.

This is the structural barrier. The whole reason APAC-Net and MFGAN cannot
solve the traffic flow MFG.

The key insight in our paper, which I will show at the end, is to sidestep
this barrier by freezing the density. Then the algebraic structure works
again. But more on that later.

% ================================================================
\section*{Slide 22 --- Part 5 separator}

\textit{(Section page.)}  Say: ``Part 5: APAC-Net.''

% ================================================================
\section*{Slide 23 --- APAC-Net}

APAC-Net was introduced by Lin and collaborators in 2021. It targets the
separable case.

APAC-Net has two networks. The first is a generator $G_\theta$, which takes
random noise $z$ and time $t$, and outputs a position $x$ in $d$-dimensional
space. As you sample $z$ many times at a fixed $t$, you get many positions,
and these positions distribute according to the population density at time
$t$. So the generator implicitly parameterises the density.

The second network is a critic $\varphi_\omega$. It takes time $t$ and
position $x$, outputs a scalar. This is the value function.

The objective comes from the Hopf formula. It is a saddle-point: the
generator minimises, the critic maximises. The training loop is the same as
a Wasserstein GAN: alternating updates of the two networks.

% ================================================================
\section*{Slide 24 --- What APAC-Net gives and what it does not}

What APAC-Net gives you, on the positive side:

You can sample the density distributionally through the generator. You get
the value function directly. You have a Wasserstein-1 convergence
guarantee. And the method scales to high dimensions.

What APAC-Net cannot do:

It cannot handle non-separable Hamiltonians. The saddle-point reformulation
breaks. In particular, it cannot solve the traffic flow MFG-LWR.

Also, no explicit convergence rate is proved. We know it converges in
Wasserstein-1, but we do not know how fast.

This is the state of the art before MFDGM. Now we are ready to read the paper.

% ================================================================
\section*{Slide 25 --- Part 6 separator}

\textit{(Section page.)}  Say: ``Part 6: the paper itself.''

% ================================================================
\section*{Slide 26 --- MFDGM: the idea}

So we come to the paper of Assouline and Missaoui from 2023.

Their key insight is very simple: stop trying to reformulate the MFG as a
min-max problem. Just plug the networks into the PDEs and minimise the
residual.

Two networks, like before. The value network $\varphi_\omega$ approximates
$\varphi$. The density network $\rho_\theta$ approximates $\rho$. Both take
time and position as inputs and output a scalar.

If the networks truly solve the MFG, then by definition the PDE residuals
are zero. So we define loss functions that measure how far the residuals
are from zero, and let an optimiser fix the weights.

The crucial point: the Hamiltonian $H$ appears directly inside the loss. We
never reformulate as a saddle-point. No Hopf formula is needed. So the
method works for any Hamiltonian, separable or not.

% ================================================================
\section*{Slide 27 --- The loss functions of MFDGM}

There are two loss functions, one per equation.

The HJB loss measures how badly the value network violates the HJB equation.
At each random collocation point in the domain, we plug $\varphi_\omega$
and $\rho_\theta$ into the HJB and square the residual. There is also a
terminal-cost penalty.

The FP loss measures how badly the density network violates the
Fokker-Planck equation. Same idea: random collocation points, plug in,
square the residual. There is also an initial-condition penalty.

The points are sampled randomly inside the domain. This is the PINN
paradigm, introduced by Raissi, Perdikaris, and Karniadakis in 2019: the
network learns the PDE solution by minimising the residual at random points.

% ================================================================
\section*{Slide 28 --- Algorithm 1 of MFDGM}

The training algorithm is exactly what you would expect.

Initialise both networks. Repeat $K$ times: sample random points, compute
the HJB loss using the current $\varphi_\omega$ and $\rho_\theta$, take one
Adam optimiser step on the value network. Then sample fresh random points,
compute the FP loss, take one Adam step on the density network.

This is alternating training. The HJB step uses the current density, and
the FP step uses the just-updated value function. Two separate Adam
optimisers, one per network.

The authors emphasise that having two separate losses with two separate
optimisers, instead of one combined loss, converges faster than DGM-MFG, an
earlier method.

% ================================================================
\section*{Slide 29 --- Part 7 separator}

\textit{(Section page.)}  Say: ``Part 7: convergence theory.''

% ================================================================
\section*{Slide 30 --- Universal Approximation, recalled}

To prove that MFDGM converges, the authors rely on the Universal
Approximation Theorem of Hornik from 1991.

The theorem says: the class of single-hidden-layer neural networks with a
$C^2$ non-constant bounded activation function is dense, in the
$C^2$-uniform sense, in continuous functions on any compact set.

In plain English: as you increase the number of neurons in the hidden
layer, you can approximate any sufficiently smooth function to arbitrary
accuracy, including its derivatives up to second order.

This is exactly what is needed for MFDGM: the network must approximate not
just the function values but also the derivatives that appear in the PDE.

% ================================================================
\section*{Slide 31 --- Theorem 3.1 (existence)}

Theorem 3.1 of the paper says: under regularity conditions H1, H2, H3, for
any positive epsilon, there exist neural networks $\rho_\theta$ and
$\varphi_\omega$ such that the residual losses are at most some constant
times epsilon.

In plain English: networks exist that make the residual arbitrarily small.

This is a statement about expressivity. It says the networks are powerful
enough. It does not yet say that we can train them; it just says good
networks exist.

The proof uses the Universal Approximation Theorem in the version for
derivatives.

% ================================================================
\section*{Slide 32 --- Theorem 3.2 (convergence)}

Theorem 3.2 is the main convergence theorem. It says, under stronger
conditions H4, H5, H6, the trained networks converge to the true MFG
solution, in $L^p$ for any $p$ less than two.

Stop and unpack this.

The trained networks converge to the true solution. Good.

In $L^p$ for $p$ less than two. This is a weak convergence statement. It is
weaker than $L^2$, much weaker than $L^\infty$, and weaker than Wasserstein.

The theorem does not give a rate. Just $\to$ as the network width goes to
infinity. No information about how fast.

So this is a positive result, but a soft one. It is the gap our work
addresses.

% ================================================================
\section*{Slide 33 --- Why L\^p for p<2 is weak}

Why exactly is $L^p$ for $p$ less than two a weak guarantee?

On a bounded domain, smaller $p$ is weaker. The chain is: $L^\infty$ is the
strongest, then $L^2$, then $L^1$. Inclusion goes from strong to weak.

To see this concretely, consider a sequence of spike functions $f_n$, equal
to $n$ on the small interval from zero to one over $n$.

The $L^1$ norm is one, constant for all $n$. The $L^2$ norm is square root
of $n$, going to infinity. The $L^\infty$ norm is $n$, going to infinity.

So a sequence can be bounded in $L^1$ but have arbitrarily tall spikes that
blow up in $L^2$. Convergence in $L^1$ is consistent with crazy spikes.
$L^2$ forbids them.

For MFDGM, $L^p$ convergence with $p$ less than two does not rule out
transient spikes in the density. It is a weak statement.

By contrast, Wasserstein-2 convergence does rule out spikes. It is
geometrically meaningful. Strengthening from $L^p$, $p$ less than two, to
Wasserstein-2 is the technical improvement our paper makes.

% ================================================================
\section*{Slide 34 --- Part 8 separator}

\textit{(Section page.)}  Say: ``Part 8: experiments.''

% ================================================================
\section*{Slide 35 --- Test 1: analytic separable benchmark}

The first experiment uses a separable MFG with a known analytic solution,
taken from the APAC-Net paper.

The Hamiltonian and the terminal cost are specifically chosen so that the
true solution can be written explicitly: $\varphi$ equals one half $x$
squared minus $d$ times $t$, and $\rho$ is a standard Gaussian.

In one dimension, the paper trains MFDGM with three hidden layers of one
hundred neurons each, using the ResNet architecture and ten thousand
training iterations.

Figure 2 of the paper shows the MFDGM solution overlaid on the analytic
solution at three time slices, $t$ equals zero point two five, zero point
five, and zero point seven five. The two curves are essentially
indistinguishable. The $L^2$ relative error drops to around one percent.

This is a sanity check, but a good one.

% ================================================================
\section*{Slide 36 --- Tests 2 to 4: scaling}

The paper does three further tests on this analytic benchmark.

Test two: vary the number of hidden units from two to fifty in a single
layer. As expected, more units gives more accuracy, shown in Figure 4.

Test three: high dimensions, $d$ equals two, fifty, and one hundred. The
residuals decay even at one hundred dimensions, shown in Figure 6.

Test four: very high dimensions, up to two hundred, with just one hidden
layer of one hundred neurons. The residuals still decay, shown in Figure 7,
and the paper later shows results at three hundred dimensions in Figure 8.

The headline of the paper: MFDGM scales to three hundred dimensions on a
single layer. This is a strong claim and it is the main technical
contribution.

% ================================================================
\section*{Slide 37 --- Comparison with APAC-Net, MFGAN, DGM-MFG}

The paper also compares MFDGM against three other deep solvers on the same
separable benchmark.

Figure 9 of the paper shows the relative error of $\rho$ and $\varphi$ for
all four methods.

MFDGM achieves the lowest error, around ten to the minus two.
MFGAN is comparable. APAC-Net is limited because its density is
distributional and has to be estimated by kernel-density methods.
DGM-MFG is slower because it uses a single combined loss.

A caveat: this comparison is on a separable benchmark, so APAC-Net and
MFGAN are allowed to play. The real test, which only MFDGM and classical
Newton can run, is the next slide.

% ================================================================
\section*{Slide 38 --- The big test: MFG-LWR traffic flow}

This is the experiment that justifies the entire paper.

The problem is the MFG-LWR traffic flow model with a non-separable
Hamiltonian. Domain is one-dimensional, time horizon one, terminal cost
zero, and the initial density is a Gaussian bump.

APAC-Net cannot run on this problem. The Hopf formula breaks. Only MFDGM
and classical Newton can.

The paper runs MFDGM both deterministically, $\nu$ equals zero, and
stochastically, $\nu$ equals zero point five. Figure 10 shows the density,
optimal cost, and speed at three different times. The solution evolves
consistently with traffic-flow intuition: density disperses, speed adjusts.

Figure 12 shows the fundamental diagram, flux versus density, which is a
classical validation in traffic theory. The MFDGM solution matches the
expected shape.

This is the unique achievement of MFDGM: a deep solver for non-separable
MFGs that scales to high dimensions.

% ================================================================
\section*{Slide 39 --- Part 9 separator}

\textit{(Section page.)}  Say: ``Part 9: our contribution. Let me describe
what our group is doing on top of this paper.''

% ================================================================
\section*{Slide 40 --- Where MFDGM leaves a gap}

Before I describe our method, let me lay out the gap MFDGM leaves open.

Look at the comparison table on the slide. Columns: handling of non-
separable Hamiltonians, representation of $\rho$, convergence metric, and
whether there is an explicit rate.

APAC-Net handles only separable, but it has distributional $\rho$,
Wasserstein-1 convergence. No rate.

MFDGM handles any Hamiltonian, the good news. But its $\rho$ is a pointwise
neural network. Its convergence is $L^p$ for $p$ less than two, which is
weak. And no explicit rate.

So neither method gives you all three properties: distributional $\rho$,
Wasserstein convergence, and an explicit rate.

The question is: can we close all three gaps simultaneously, at least in
some useful regime, without giving up non-separable handling?

% ================================================================
\section*{Slide 41 --- Our insight: freeze the density}

The key insight is very simple. Here it is.

If $\rho$ is held fixed, any non-separable Hamiltonian becomes separable.

Look at the diagram. The full non-separable MFG on the left. We freeze the
density at the current iterate $\rho^k$. Then the effective Hamiltonian
$H^k$ of $x$, $p$ is just $H_0$ plus epsilon times $R$ evaluated at $\rho^k$.
Critically, this depends only on $x$ and $p$. It is separable in the sense
the GAN solvers require.

So we can apply APAC-Net to this separable subproblem. APAC-Net gives back
a new density $\rho^{k+1}$. We update and repeat.

This is essentially a fixed-point iteration on the density. Each step is one
APAC-Net call.

The whole research question is: when does this outer iteration converge,
and how fast?

% ================================================================
\section*{Slide 42 --- Weakly non-separable Hamiltonians}

To answer the question, we define the right class of Hamiltonians.

We say $H$ is weakly non-separable with strength epsilon if it splits into
three parts: a strongly convex base $H_0$ with convexity constant $c_0$;
a Lasry-Lions monotone interaction $f_0$ with constant lambda; and a
coupling residual $R$, scaled by epsilon, that is Lipschitz in $\rho$ with
constant $L_R$ and Lipschitz in $p$ with constant $L_{R,p}$.

Special cases. When epsilon is zero, you recover the standard separable
MFG. When $R$ is rho times $p$, you recover the traffic flow model. When
$R$ depends only on density, $L_{R,p}$ is zero.

The most important quantity is the boxed condition at the bottom: epsilon
times $L_R$ strictly less than lambda. We call this the GAN-tractable
regime. The threshold balances coupling strength against the restoring
force from monotonicity.

Below the threshold, our method works. Above the threshold, it fails.

% ================================================================
\section*{Slide 43 --- The Extended-MFGAN algorithm}

The algorithm is the obvious thing.

Initialise the density. Then loop. At each step, build the effective
Hamiltonian by plugging in the current density. Run an inner GAN solver, for
example APAC-Net, on the resulting separable MFG. Update the density.
Repeat.

A nice sanity check: when epsilon is zero, the effective Hamiltonian does
not depend on the current density. The outer loop converges in one step,
and Extended-MFGAN reduces exactly to APAC-Net.

The total cost is $K$ outer iterations times one APAC-Net call. Once we
prove the contraction rate kappa, we have an explicit formula for how
many outer iterations are needed: log one over delta divided by log one
over kappa, for any desired accuracy delta.

% ================================================================
\section*{Slide 44 --- Our main theorem: L\^2 contraction}

Our main result is the boxed theorem on the slide.

Under monotonicity, strong convexity, bounded density, and the GAN-tractable
condition, the outer iteration is an $L^2$ contraction. The error after
step $k$ plus one is at most kappa times the error after step $k$, where
kappa equals epsilon times $L_R$ divided by lambda.

This implies Wasserstein-2 convergence at rate kappa to the $k$ over two.
The square-root scaling comes from the standard $L^2$-to-Wasserstein
inequality.

This is the first explicit convergence rate for a deep MFG solver in the
non-separable setting.

The proof is via the Lasry-Lions duality identity, the same technique used
to prove uniqueness of the MFG equilibrium in the classical theory.

We also prove the rate is sharp on the bilinear traffic-flow family: the
linearised outer map has operator norm exactly epsilon times absolute value
of one minus epsilon divided by lambda.

And we prove the condition is necessary: when epsilon $L_R$ is at least
lambda, the iteration diverges, by an explicit linearised counterexample.

% ================================================================
\section*{Slide 45 --- Comparison: where we sit}

Let me put our method back in the comparison table.

APAC-Net: separable only, distributional, Wasserstein-1, no rate.

MFDGM: non-separable, pointwise, $L^p$ for $p$ less than two, no rate.

Our method, Extended-MFGAN: handles the weakly non-separable regime,
distributional, $L^2$ and Wasserstein-2 convergence, explicit rate kappa
to the $k$.

We close all three gaps simultaneously, within the weakly non-separable
regime.

The trade-off is honest. We are restricted to epsilon $L_R$ less than
lambda. Above this threshold, MFDGM remains the right tool. Our method
complements MFDGM rather than competing with it.

% ================================================================
\section*{Slide 46 --- Summary}

Let me summarise everything.

The paper of Assouline and Missaoui. They introduce MFDGM, a physics-
informed deep solver for Mean Field Games. It uses two neural networks
trained alternately to minimise the HJB and FP residuals. It works for any
Hamiltonian, including non-separable ones like traffic flow. It scales to
three hundred dimensions on a single hidden layer. The convergence
guarantee is $L^p$ for $p$ less than two, with no explicit rate.

Our extension. We introduce weakly non-separable Hamiltonians, a controlled
interpolation between separable and non-separable. We propose a fixed-point
algorithm called Extended-MFGAN: freeze the density inside the residual,
solve the resulting separable subproblem with an existing GAN solver,
update, and repeat. We prove $L^2$ contraction with explicit rate epsilon
$L_R$ over lambda, which gives Wasserstein-2 convergence. The threshold
is sharp: we prove necessity by an explicit counterexample.

Importantly, the two methods are complementary, not competing. MFDGM is the
right tool above the threshold. Extended-MFGAN is the right tool below the
threshold.

% ================================================================
\section*{Slide 47 --- Thank you}

This brings me to the end of the talk.

In one sentence: I walked us through a recent paper that pushes deep
learning for Mean Field Games into the non-separable regime, and I showed
how our follow-up work strengthens its guarantees to Wasserstein-2 with
an explicit convergence rate.

Thank you for your attention. I am very happy to take questions.

\end{document}
